\documentclass[11pt,a4paper,fleqn]{article}
%\geometry{landscape}                % Activate for for rotated page geometry
%\usepackage[parfill]{parskip}    % Activate to begin paragraphs with an empty line rather than an indent
\usepackage{graphicx}
\usepackage{amssymb}
\usepackage{amsmath}
\usepackage{epstopdf}
\include{rvc-notation}
\usepackage{matlab}
\DeclareGraphicsRule{.tif}{png}{.png}{`convert #1 `dirname #1`/`basename #1 .tif`.png}

\title{Solving trigonometric equations}
\author{Peter Corke}
%\date{}                                           % Activate to display a given date or no date
\setlength{\parindent}{0pt}

\begin{document}
\maketitle
%\section{}
%\subsection{}
When solving problems in inverse kinematics we often end up with an equation of the form.  There are several ways to solve this and they
are discussed below.

\section*{Simple derivation}
\begin{equation}
a \cos \theta + b \sin \theta = c, \,\, a,b \ne 0 \label{eq:problem-orig}
\end{equation}
and we note similarity with the form of the sum of angles identity
\[
 \underbrace{\sin \phi}_{\sim a} \cos \theta + \underbrace{\cos \phi}_{\sim b} \sin \theta \equiv \sin (\phi + \theta) 
\]
In order to equate coefficients we need to ensure that $\sin^2 \phi + \cos^2 \phi = 1$ which is true only if $a^2 + b^2 = 1$. 
In general this will not be the case so we normalize the equation, dividing each side by $d = \sqrt{a^2 + b^2}$ giving
\begin{equation}
a^\prime \cos \theta + b^\prime \sin \theta = c^\prime \label{eq:problem}
\end{equation}
where $a^\prime= a/d$, $b^\prime=b/d$ and $c^\prime=c/d$.
Now we can write
\[
\sin \phi = a^\prime, \, \cos \phi = b^\prime
\]
and solve for $\phi$
\[
\tan \phi = \frac{a^\prime}{b^\prime}  = \frac{a}{b} \in [-2\pi, 2\pi)
\]
which should be computed using an \texttt{atan2} function. 

Next we  rewrite \eq{eq:problem} as
\[
\sin (\phi + \theta) = c^\prime
\]
and if $|c^\prime| \le 1$ or $a^2+b^2-c^2 > 0$ we can solve  for
\[
\theta = \sin^{-1} c^\prime - \phi
\]

In general, there is a second solution corresponding to the negative solution of the square root $d = -\sqrt{a^2 + b^2}$ leading to
\[
\tan \phi = \frac{-a^\prime}{-b^\prime}  = \frac{-a}{-b} \in [-2\pi, 2\pi)
\]
which puts the solution for $\phi$ in the diagonally opposite quadrant, and
\[
\sin (\phi + \theta) = -c^\prime
\]
Since $\sin( -x) = -\sin(x)$ we can write the second solution as
\[
\theta = -\sin^{-1} c^\prime - \phi
\]

In summary, the two solutions are
\[
\begin{array}{ll}
\theta = \hspace{1em}\sin^{-1} c^\prime - \phi, & \tan \phi =  \frac{a}{b} \\
\theta = -\sin^{-1} c^\prime - \phi, & \tan \phi =  \frac{-a}{-b}
\end{array}
\]

We can test this numerically using MATLAB
\begin{matlabcode}
a = 4; b = 5; c = 3;
clear theta

phi = atan2(a, b);
th1 = asin(c/norm([a b])) - phi
\end{matlabcode}
\begin{matlaboutput}
th1 = -0.1871
\end{matlaboutput}
\begin{matlabcode}
phi = atan2(-a, -b);
th2 = -asin(c/norm([a b])) - phi;
theta = [th1 th2];
a*cos(theta) + b*sin(theta) - c
\end{matlabcode}
\begin{matlaboutput}
ans = 1x2    
1.0e+-15 *
   -0.8882   -0.8882
\end{matlaboutput}
which indicates solutions equal to zero up to machine precision.

\section*{Other forms}
Another commonly given solutions of this equation include
\[
\theta = \tan^{-1} \frac{c}{\pm \sqrt{a^2+b^2-c^2}} - \tan^{-1}\frac{a}{b}
\]
and
\[
\theta = \tan^{-1} \frac{ bc \pm ad}{ac \mp bd}
\]
which requires just a single arc-tangent operation.
See the discussion at \url{https://math.stackexchange.com/questions/213545/solving-trigonometric-equations-of-the-form-a-sin-x-b-cos-x-c}


\section*{The Weierstrass transformation}
A well known way to convert trigonometric equations to algebraic equations is with the Weierstrass transformation\footnote{After the German mathematician Karl Weierstrass (1815-1897).} which is familiar as one of the half-angle identities
\[
\sin \theta = \frac{2h}{1+h^2}, \,\,\, \cos \theta = \frac{1-h^2}{1+h^2}
\]
where $h = \tan \frac{\theta}{2}$.  The problem (\ref{eq:problem-orig}) can be rewritten as
\[
\frac{2 b h}{h^2 + 1} - \frac{a (h^2 - 1)}{h^2 + 1} = c
\]
which can be expressed as a quadratic in $h$
\[
(a+c)h^2 - 2bh+a + c = 0
\]
which we can solve as
\begin{align}
h = \left(\begin{array}{c} \frac{b+\sqrt{a^2+b^2-c^2}}{a+c}\\ 
\frac{b-\sqrt{a^2+b^2-c^2}}{a+c} \end{array}\right)
\end{align}
and clearly has the two solutions so long as $a^2 + b^2 - c^2 > 0$. 
The transformation has led to the solution in a very straightforward fashion.

Once again, we can test this numerically using MATLAB
\begin{matlabcode}
syms a b c h theta
e = a*cos(theta) + b*sin(theta) == c
\end{matlabcode}
\begin{matlabsymbolicoutput}
e = 
    $\displaystyle a \cos \left(\theta \right)+b \sin \left(\theta \right)=c$
\end{matlabsymbolicoutput}
\begin{matlabcode}
e = rewrite(e, 'tan')
\end{matlabcode}
\begin{matlabsymbolicoutput}
e = 
    $\displaystyle \frac{2 b \tan \left(\frac{\theta }{2}\right)}{{\tan \left(\frac{\theta }{2}\right)}^2 +1}-\frac{a {\left({\tan \left(\frac{\theta }{2}\right)}^2 -1\right)}}{{\tan \left(\frac{\theta }{2}\right)}^2 +1}=c$
\end{matlabsymbolicoutput}
\begin{matlabcode}
e = subs(e, tan(theta/2), h)
\end{matlabcode}
\begin{matlabsymbolicoutput}
e = 
    $\displaystyle \frac{2 b h}{h^2 +1}-\frac{a {\left(h^2 -1\right)}}{h^2 +1}=c$
\end{matlabsymbolicoutput}
\begin{matlabcode}
sol = solve(e, h)
\end{matlabcode}
\begin{matlabsymbolicoutput}
sol = 
    $\displaystyle \left(\begin{array}{c}
\frac{b+\sqrt{a^2 +b^2 -c^2 }}{a+c}\\
\frac{b-\sqrt{a^2 +b^2 -c^2 }}{a+c}
\end{array}\right)$
\end{matlabsymbolicoutput}

\begin{par}
\begin{flushleft}
Now lets validate this
\end{flushleft}
\end{par}

\begin{matlabcode}
a = 4; b = 5; c = 3;
theta = 2*atan(eval(sol))
\end{matlabcode}
\begin{matlaboutput}
theta = 2x1    
    1.9792
   -0.1871

\end{matlaboutput}
\begin{matlabcode}
a*cos(theta) + b*sin(theta) - c
\end{matlabcode}
\begin{matlaboutput}
ans = 2x1    
1.0e+-15 *
         0
   -0.8882
\end{matlaboutput}
which again is equal to zero up to machine precision.


\end{document}  